% texforge example — a technical report/article.
% Renders with: texforge render examples/report/main.tex
\documentclass[11pt,a4paper]{article}

\usepackage[margin=2.5cm]{geometry}
\usepackage{microtype}          % better justification; fixes many overfull boxes
\usepackage{amsmath,amssymb}
\usepackage{booktabs}
\usepackage{graphicx}
\usepackage{xcolor}
\usepackage{hyperref}
\hypersetup{colorlinks=true, linkcolor=blue!60!black, urlcolor=teal!70!black}

\title{\textbf{texforge}\\\large A self-healing \LaTeX\,$\rightarrow$\,PDF engine}
\author{Rendered by the texforge engine on tectonic}
\date{\today}

\begin{document}
\maketitle

\begin{abstract}
This report is a fixture: it exercises sectioning, cross-references, a table, an
equation and a bulleted list so a single build proves the toolchain end to end.
Because tectonic fetches missing packages on demand, this file compiles from a
clean machine with no \TeX{} installation.
\end{abstract}

\tableofcontents
\vspace{1em}

\section{Introduction}
\label{sec:intro}
texforge wraps the tectonic engine with an observable CLI. It streams structured
logs, times every build, and reads the engine's own output to suggest fixes for
common problems --- see Section~\ref{sec:results}.

\section{A little mathematics}
The Gaussian integral, for good measure:
\begin{equation}
  \int_{-\infty}^{\infty} e^{-x^2}\,\mathrm{d}x = \sqrt{\pi}.
  \label{eq:gauss}
\end{equation}
Equation~\eqref{eq:gauss} typesets cleanly with \texttt{amsmath}.

\section{Results}
\label{sec:results}
A small table via \texttt{booktabs}:

\begin{table}[h]
  \centering
  \begin{tabular}{lrr}
    \toprule
    Stage & Cold build & Warm build \\
    \midrule
    Fetch packages & yes & cached \\
    Compile passes & 2   & 2 \\
    Total          & slower & fast \\
    \bottomrule
  \end{tabular}
  \caption{Cold vs.\ warm builds. Warm builds reuse the package cache.}
  \label{tab:builds}
\end{table}

Design goals:
\begin{itemize}
  \item \textbf{Self-healing} --- packages arrive on demand.
  \item \textbf{Self-optimizing} --- caching plus log-driven suggestions.
  \item \textbf{Observable} --- level-controlled, structured logs.
  \item \textbf{Portable} --- one static binary per OS, no runtime deps.
\end{itemize}

\section{Conclusion}
If you can read this PDF, the engine works. See the project
\href{https://github.com/m3di/texforge}{README} for the quickstart.

\end{document}
